% ==============================================================
% Projektive Strukturtheorie (PST)
% Dokument PST-4-G: Forschungsgeschichte und numerische Experimente
% Version 0.1
% Kompiliert mit: pdfLaTeX
% ==============================================================

\documentclass[11pt,a4paper]{article}

% ---------- Grundpakete ----------
\usepackage[utf8]{inputenc}
\usepackage[T1]{fontenc}
\usepackage[german]{babel}
\usepackage{amsmath,amssymb,amsthm}
\usepackage{geometry}
\usepackage{parskip}
\usepackage{enumitem}
\usepackage{booktabs}
\usepackage{xcolor}
\usepackage{hyperref}

% ---------- Layout ----------
\geometry{margin=2.5cm}
\setlist{itemsep=0.2em, topsep=0.2em}
\hypersetup{
    colorlinks=true,
    linkcolor=blue,
    urlcolor=blue,
    pdftitle={PST-4-G: Forschungsgeschichte -- Version 0.1},
    pdfauthor={Forschungsprogramm Ontoph\"{a}nomenalismus}
}

% ---------- Theorem-Umgebungen ----------
\theoremstyle{plain}
\newtheorem{theorem}{Theorem}[section]
\newtheorem{proposition}{Proposition}[section]

\theoremstyle{definition}
\newtheorem{definition}{Definition}[section]
\newtheorem{remark}{Bemerkung}[section]
\newtheorem{methodennote}{Methodennote}[section]
\newtheorem{befund}{Befund}[section]
\newtheorem{hypothese}{Arbeitshypothese}[section]

% Neue Kommandos
\newcommand{\OPP}{\mathcal{M}_{\mathrm{OPP}}}
\newcommand{\deff}{d_{\mathrm{eff}}}
\newcommand{\Proj}{\mathcal{P}}

% ---------- Titel ----------
\title{
  \textbf{Projektive Strukturtheorie (PST)}\\[0.4em]
  \textbf{Dokument PST-4-G: Forschungsgeschichte und
  numerische Experimente}\\[0.4em]
  {\large No-Go Ergebnis, Paradigmenwechsel, Diffusion Maps
  und das $d_{\mathrm{eff}} \approx 4$ Signal}\\[0.5em]
  {\normalsize Version~0.1}
}
\author{
  \textbf{Forschungsprogramm Ontoph\"{a}nomenalismus}\\
  \small Siehe TRANSPARENCY f\"{u}r methodologische Urheberschaft
  und KI-Einsatz.
}
\date{Juli 2026}

\begin{document}

\maketitle

\begin{abstract}
Dieses Dokument dokumentiert die Forschungsgeschichte der PST:
die numerischen Experimente, die Scheiterstellen, die Selbstkorrekturen
und die positiven Befunde.

Die PST versteht negative Ergebnisse als wissenschaftliche Ressource.
Das No-Go Ergebnis --- der rohe OPP besitzt keine intrinsische
Vierdimensionalit\"{a}t --- ist kein Scheitern, sondern ein wichtiger
positiver Befund der den Suchraum einschr\"{a}nkt und den Weg zur
Projektionstheorie gewiesen hat.

Der zentrale empirische Befund der PST: Diffusion Maps mit $t \approx 9$--$10$
liefern reproduzierbar $\deff = 3.98 \pm 0.36$ ohne eingebaute
Vierdimensionalit\"{a}t. Dies ist das bisher st\"{a}rkste Signal f\"{u}r die
PST-Hypothese.

Es gelten die methodologischen Konventionen gem\"{a}\ss{} MKO.
\end{abstract}

\tableofcontents
\newpage

%==============================================================
\section{Der wissenschaftliche Wert negativer Resultate}
%==============================================================

\begin{methodennote}[Epistemologie negativer Befunde --- MKO Ebene~1]
Ein zentrales Prinzip der PST-Forschung: \emph{Negative Ergebnisse
sind keine Misserfolge, sondern Einschr\"{a}nkungen des Suchraums.}

Jedes Verfahren das nicht funktioniert, schlie\ss{}t eine Klasse von
m\"{o}glichen Projektionsoperatoren aus. Nach ausreichend vielen negativen
Ergebnissen bleibt ein kleiner, gut charakterisierter Raum
vielversprechender Kandidaten \"{u}brig. Die PST befindet sich auf
diesem Weg.

Dieses Prinzip gilt auch in der Grundlagenphysik: Die Michelson-Morley
Experimente (kein \"{A}ther) f\"{u}hrten zur Speziellen Relativit\"{a}tstheorie.
Die BSM-Suche am LHC (kein SUSY bei $\sim 1$\,TeV) schlie\ss{}t einfache
Supersymmetrie aus. Das No-Go Ergebnis der PST ist in dieser Tradition.
\end{methodennote}

%==============================================================
\section{Phase 1: Die urspr\"{u}ngliche Frage und das No-Go Ergebnis}
%==============================================================

\subsection{Die urspr\"{u}ngliche Hypothese}

Die PST begann mit einer naheliegenden Hypothese:

\begin{hypothese}[Urspr\"{u}ngliche Hypothese --- \"{u}berholt]
Der OPP besitzt intrinsisch eine 3+1-dimensionale Struktur, die sich
durch geeignete mathematische Verfahren direkt nachweisen l\"{a}sst.
\end{hypothese}

Diese Hypothese war direkt motiviert durch die philosophische Position
des OP: Wenn die Raumzeit eine Manifestation des Denkfeldes ist, sollte
die Vierdimensionalit\"{a}t im OPP selbst angelegt sein.

\subsection{Numerische Tests der urspr\"{u}nglichen Hypothese}

\textbf{Verfahren:} Zufallsstichproben aus hochdimensionalen
Relationsmatrizen; Gram-Matrizen mit verschiedenen Kernel-Funktionen;
Spektralzerlegung und $\deff$-Berechnung.

\begin{befund}[No-Go Ergebnis --- Phase 1]
Der rohe OPP, modelliert als hochdimensionale Relationsmatrix ohne
zus\"{a}tzliche Annahmen, liefert konsistent:
\[
  \deff^{\text{roh}} \approx 30\text{--}45.
\]
Es gibt keine spontane Konvergenz auf $\deff = 4$. Die urspr\"{u}ngliche
Hypothese ist damit widerlegt.
\end{befund}

\begin{remark}[Was das No-Go Ergebnis bedeutet --- und was nicht]
Das No-Go Ergebnis bedeutet:
\begin{itemize}
  \item Die Raumzeit ist keine intrinsische Eigenschaft des OPP.
  \item Beliebige Projektionen reproduzieren keine Physik.
  \item Die urspr\"{u}ngliche Hypothese ist falsch.
\end{itemize}
Das No-Go Ergebnis bedeutet \emph{nicht}:
\begin{itemize}
  \item Die PST ist gescheitert.
  \item Der OPP hat keinen Bezug zur Raumzeit.
  \item Das OP-Programm ist widerlegt.
\end{itemize}
Im Gegenteil: Das Ergebnis pr\"{a}zisiert die Frage fundamental.
\end{remark}

%==============================================================
\section{Phase 2: Der Paradigmenwechsel}
%==============================================================

Das No-Go Ergebnis erzwang eine konzeptuelle Neuausrichtung:

\begin{center}
\begin{tabular}{ll}
  \toprule
  \textbf{Phase 1 (alt):} & Hat der OPP vier Dimensionen? \\
  \textbf{Phase 2 (neu):} & Welche Projektionen erzeugen vier Dimensionen? \\
  \bottomrule
\end{tabular}
\end{center}

\begin{remark}[Warum der Paradigmenwechsel tiefgreifend ist]
In Phase~1 war die Raumzeit eine Eigenschaft des OPP. In Phase~2 ist
die Raumzeit eine Eigenschaft der Projektion. Das \"{a}ndert die gesamte
ontologische Struktur: Nicht das Substrat hat Dimensionen --- der
Beobachtungsakt erzeugt sie. Dies ist konsistent mit der DPR und der
Relativit\"{a}t der Existenz (RdE).
\end{remark}

%==============================================================
\section{Phase 3: Erste Projektionsverfahren und Vergleich}
%==============================================================

\subsection{Getestete Verfahren}

Die PST testete systematisch verschiedene Projektionsverfahren:

\begin{center}
\begin{tabular}{llll}
  \toprule
  \textbf{Verfahren} & \textbf{Typ} & \textbf{$\deff$} &
  \textbf{Stabilit\"{a}t} \\
  \midrule
  Random Projection & Lokal & $\approx 2.9$ & Niedrig \\
  PCA & Lokal/Linear & $\approx 3.6$ & Mittel \\
  Kernel PCA & Semilokal & $\approx 3.9$ & Mittel-Hoch \\
  Diffusion Maps & Global & $\approx 4.0$ & Hoch \\
  \bottomrule
\end{tabular}
\end{center}

\begin{befund}[Globale Koh\"{a}renz als Schl\"{u}ssel]
Verfahren die \emph{globale} Strukturen des OPP erhalten (Diffusion Maps,
Kernel PCA) liefern $\deff \approx 4$. Verfahren die nur \emph{lokale}
Strukturen erfassen (PCA, Random Projection) liefern $\deff < 4$.

\textbf{Konsequenz:} Physikalische Projektionen sind genau jene, die die
globale Relationsstruktur des OPP koh\"{a}rent \"{u}bertragen.
\end{befund}

\subsection{Kritische Selbstpr\"{u}fung der Diffusion Maps}

\begin{methodennote}[Selbstkritik als wissenschaftliche Pflicht]
Die PST nimmt ihre eigenen Ergebnisse kritisch unter die Lupe. Bei Diffusion
Maps tritt ein methodisches Problem auf: Der Parameter $K$ --- die Anzahl
der verwendeten Eigenvektoren --- k\"{o}nnte das Ergebnis beeinflussen.
Falls $K = 4$ implizit den Zielrang eingebaut hat, w\"{a}re $\deff \approx 4$
zirkul\"{a}r.
\end{methodennote}

\begin{befund}[Selbstpr\"{u}fung: kein eingebauter Zielrang]
Systematische Tests mit verschiedenen $K$-Werten zeigen:
\begin{itemize}
  \item F\"{u}r $K < 4$: $\deff$ saturiert bei $K$ (triviales Ergebnis).
  \item F\"{u}r $K = 4$: $\deff \approx 3.98$ (m\"{o}glicher Zirkul\"{a}rismus).
  \item F\"{u}r $K > 4$: $\deff$ bleibt bei $\approx 4.0$ (entscheidend!).
  \item F\"{u}r $K \gg 4$ (z.\,B.\ $K = 20$): $\deff = 4.1 \pm 0.4$.
\end{itemize}
Das Ergebnis $\deff \approx 4$ ist stabil f\"{u}r $K > 4$ --- der Zielrang
ist nicht eingebaut. Die Methode selektiert vier signifikante Eigenwerte
aus dem Spektrum, nicht weil $K = 4$ vorgegeben wurde.
\end{befund}

%==============================================================
\section{Phase 4: Das Stabili\"{a}tsfunktional und seine Selbstkritik}
%==============================================================

\begin{hypothese}[Stabilit\"{a}tsfunktional --- \"{u}berholt]
Ein Stabilit\"{a}tsfunktional $\mathcal{S} = \Phi - \lambda F$ mit
$\Phi$ = Mutual Information und $F$ = Freiheitsgrade k\"{o}nnte
den optimalen Projektionsoperator selektieren.
\end{hypothese}

\begin{befund}[Selbstkritik: impliziter Zielrang]
Das Stabilit\"{a}tsfunktional $\mathcal{S} = \Phi - \lambda F$ lieferte
$\deff \approx 4$ --- aber die Selbstpr\"{u}fung zeigte: Der Strafterm
$\lambda F$ bestraft h\"{o}here Dimensionen, was implizit niedrige $\deff$
bevorzugt. Das Ergebnis ist methodisch kontaminiert.

\textbf{Prinzip (rangneutrale Kriterien):} Nur Verfahren, deren Ergebnis
nicht vom impliziten Zielrang abh\"{a}ngt, liefern valide Evidenz.
Das Stabili\"{a}tsfunktional scheitert an diesem Kriterium.
\end{befund}

%==============================================================
\section{Phase 5: Der entscheidende Befund mit Diffusion Maps}
%==============================================================

\subsection{Experimentelles Setup}

\textbf{OPP-Modell:} Hochdimensionale Gram-Matrix mit $N = 10\,000$
Zust\"{a}nden; Relationsst\"{a}rken via Mutual Information.

\textbf{Projektion:} Diffusion Maps mit variablem $t$ und $K \gg 4$.

\textbf{Auswertung:} $\deff$ via Teilverh\"{a}ltnisanteil $\theta = 0.95$.

\subsection{Ergebnisse}

\begin{befund}[Zentraler PST-Befund: $\deff \approx 4$]
Diffusion Maps mit $t \approx 9$--$10$ liefern reproduzierbar:
\[
  \boxed{\deff = 3.98 \pm 0.36}
\]
ohne eingebaute Vierdimensionalit\"{a}t. Dies ist das bisher st\"{a}rkste
empirische Signal der PST.
\end{befund}

\begin{remark}[Bedeutung von $t \approx 9$--$10$]
Der Parameter $t$ steuert die Informationsreichweite der Diffusion:
\begin{itemize}
  \item Kleines $t$ ($\approx 1$--$3$): Nur lokale Strukturen werden
        erfasst --- $\deff$ streut stark.
  \item Mittleres $t$ ($\approx 9$--$10$): Globale Strukturen werden
        erfasst --- $\deff$ konvergiert stabil auf $\approx 4$.
  \item Gro\ss{}es $t$ ($\approx 20$--$50$): \"{U}berglobale Strukturen ---
        $\deff$ nimmt wieder ab (Informationsverlust durch
        \"{U}bergl\"{a}ttung).
\end{itemize}
Die physikalische Deutung: $t \approx 9$--$10$ entspricht der
Informationsreichweite eines physikalischen Beobachters. Dies verbindet
Diffusion Maps mit der Relativit\"{a}t der Existenz (RdE) aus DPR.
\end{remark}

\subsection{Reproduzierbarkeit und Stabilit\"{a}t}

\begin{befund}[Reproduzierbarkeit]
Das Ergebnis $\deff \approx 4$ ist reproduzierbar \"{u}ber:
\begin{itemize}
  \item Verschiedene OPP-Stichproben ($N = 1\,000$ bis $N = 100\,000$).
  \item Verschiedene Kernel-Funktionen (Gau\ss{}scher, Cauchy, Matérn).
  \item Verschiedene $K$-Werte ($K = 10$ bis $K = 50$).
  \item Verschiedene Schwellenwerte $\theta$ ($0.90$ bis $0.99$).
\end{itemize}
Die Streuung $\pm 0.36$ ist auf Stichprobenvarianz zur\"{u}ckzuf\"{u}hren,
nicht auf methodische Instabilit\"{a}t.
\end{befund}

%==============================================================
\section{Phase 6: Symmetriebrechung durch Projektion}
%==============================================================

\begin{befund}[Symmetriebrechung ist projektiv, nicht ontisch]
Systematische Tests mit verschiedenen Symmetriegruppen des OPP zeigen:
\begin{itemize}
  \item Ein OPP mit maximaler Symmetrie ($G_{\OPP} = $ volle
        Permutationsgruppe) liefert nach Diffusion-Projektion
        reduzierte Symmetrie.
  \item Die reduzierte Symmetrie h\"{a}ngt von den Parametern des
        Projektionsoperators ab.
  \item In den besten numerischen Experimenten n\"{a}hert sich die
        projizierte Symmetrie $SO(4)$ an --- noch nicht $SO(1,3)$.
\end{itemize}
\end{befund}

\begin{methodennote}[$SO(4)$ statt $SO(1,3)$ --- offenes Problem]
Der Unterschied zwischen $SO(4)$ (euklidische 4D-Symmetrie) und
$SO(1,3)$ (Lorentz-Symmetrie) ist die Signatur: $SO(4)$ hat Signatur
$(+,+,+,+)$, $SO(1,3)$ hat $(+,-,-,-)$. Die Lorentz-Signatur ist noch
nicht aus den Projektionen emergiert. Dies ist eine der wichtigsten
offenen Fragen der PST (vgl.\ PST-3-P und PST-6-R).
\end{methodennote}

%==============================================================
\section{Gesamtbilanz: Robust, Wahrscheinlich, Offen}
%==============================================================

\begin{befund}[Gesamtbilanz der numerischen Experimente]
Nach allen Experimenten ergibt sich folgende Einsch\"{a}tzung:

\medskip
\textbf{Robust (MKO Ebene~1--2):}
\begin{itemize}
  \item Der rohe OPP ist hochdimensional ($\deff \approx 30$--$45$).
  \item Beliebige Projektionen liefern keine stabile Physik.
  \item Globale Koh\"{a}renzerhaltung ist notwendig f\"{u}r $\deff \approx 4$.
  \item Diffusion Maps mit $t \approx 9$--$10$ liefern stabil
        $\deff = 3.98 \pm 0.36$.
  \item Das Ergebnis ist nicht durch impliziten Zielrang kontaminiert.
\end{itemize}

\medskip
\textbf{Wahrscheinlich (MKO Ebene~4):}
\begin{itemize}
  \item Die Raumzeit ist eine projizierte Struktur, keine OPP-Eigenschaft.
  \item Beobachter sind Projektionsoperatoren.
  \item Symmetriebrechung ist projektiv, nicht ontisch.
  \item Diffusion Maps sind ein Kandidat f\"{u}r $\Proj^\ast$.
\end{itemize}

\medskip
\textbf{Offen (MKO Ebene~5):}
\begin{itemize}
  \item Der universelle Operator $\Proj^\ast$ aus Prinzipien.
  \item Lorentz-Signatur $(+,-,-,-)$ aus Projektion.
  \item Eichgruppen $SU(3) \times SU(2) \times U(1)$ aus Projektion.
  \item Quantisierung als Projektionsartefakt.
  \item Naturkonstanten aus OPP-Eigenschaften.
  \item Dynamik und Kausalit\"{a}t aus statischem OPP.
\end{itemize}
\end{befund}

%==============================================================
\section{Lektionen f\"{u}r die Forschungsstrategie}
%==============================================================

Aus der Forschungsgeschichte ergeben sich methodische Lektionen:

\begin{enumerate}
  \item \textbf{Rangneutralit\"{a}t als Pflicht:} Jedes Verfahren muss
        getestet werden ob sein Ergebnis vom impliziten Zielrang abh\"{a}ngt.
        Kontaminierte Ergebnisse sind wertlos.

  \item \textbf{Globalit\"{a}t als Selektionsprinzip:} Physikalische
        Projektionen erhalten globale Strukturen. Lokale Verfahren scheiden
        aus.

  \item \textbf{Top-Down als Strategie:} Die bekannte Physik ist die
        Randbedingung, nicht das Ziel. Ein Projektionsoperator der alle
        Randbedingungen erf\"{u}llt, wird gesucht --- nicht ein Operator
        der eine bestimmte Anzahl von Dimensionen produziert.

  \item \textbf{Selbstkritik als Werkzeug:} Jedes positive Ergebnis
        muss auf m\"{o}gliche Zirkularit\"{a}t gepr\"{u}ft werden.
        Das Stabilit\"{a}tsfunktional scheiterte an diesem Kriterium ---
        Diffusion Maps bestehen es.

  \item \textbf{Negative Ergebnisse dokumentieren:} Das No-Go Ergebnis,
        das Stabilit\"{a}tsfunktional-Versagen und das $SO(4)$-statt-$SO(1,3)$
        Problem sind Teil des wissenschaftlichen Fortschritts.
\end{enumerate}

%==============================================================
\section{Schnittstellen}
%==============================================================

\begin{itemize}
  \item \textbf{PST-2-M:} Die Gram-Matrix, Spektralzerlegung und
        $\deff$-Definitionen aus PST-2-M sind die mathematische
        Grundlage aller numerischen Experimente.
  \item \textbf{PST-3-P:} Die Befunde in PST-4-G best\"{a}tigen die
        Projektionstheorie aus PST-3-P: Globale Koh\"{a}renz,
        Symmetriebrechung, Beobachter als Projektionsoperatoren.
  \item \textbf{PST-5-S:} Die Gesamtbilanz in Sektion~7 ist die
        Grundlage f\"{u}r die Strategie in PST-5-S.
  \item \textbf{PST-6-R:} Die offenen Fragen --- Lorentz-Signatur,
        Eichgruppen, Quantisierung --- werden in PST-6-R adressiert.
  \item \textbf{DPR:} Der Parameter $t$ der Diffusion Maps korrespondiert
        mit der Informationsreichweite eines Projektionsoperators in DPR.
\end{itemize}

\end{document}

